\section{Affine rays and the direct-entry sign}
\label{sec:rays-signs}

\subsection{Exact ray spacing}

\begin{theorem}[General fixed-\((q,u,v)\) direct-entry ray]
\label{thm:direct-entry-ray}
Fix \(m,n,h_0\), and suppose one actual key realizes coordinates
\((q,u,v)\), with data
\[
 (j_0,Q_{m,0},Q_{n,0})
\]
exists, define
\begin{equation}
 \lambda_{mn}(q,u,v)
 =
 \operatorname{lcm}\left(
       \frac{qu}{(m,qu)},
       \frac{qv}{(n,qv)}
      \right)
 \label{eq:general-ray-spacing}
\end{equation}
Then every other physical key with the same \((q,u,v)\) has a unique
integer \(z\) for which
\begin{align}
 j&=j_0+\lambda_{mn}(q,u,v)\,z,
 \label{eq:direct-time-ray}\\
 Q_m&=Q_{m,0}+\frac{m\lambda_{mn}(q,u,v)}{qu}\,z,
 \label{eq:direct-Qm-ray}\\
 Q_n&=Q_{n,0}+\frac{n\lambda_{mn}(q,u,v)}{qv}\,z
 \label{eq:direct-Qn-ray}
\end{align}
\end{theorem}

\begin{proof}
For two keys with the same coordinates, take differences.  From
\[
 mj+h_0=quQ_m,\qquad nj+h_0=qvQ_n
\]
one obtains
\[
 m\,\delta j=qu\,\delta Q_m,\qquad
 n\,\delta j=qv\,\delta Q_n.
\]
The first divisibility is equivalent to
\(qu/(m,qu)\mid\delta j\), and the second to
\(qv/(n,qv)\mid\delta j\).  Their simultaneous solution lattice is
therefore exactly
\(\delta j=\lambda_{mn}(q,u,v)z\).  Substitution gives the two displayed
prime increments; they are integers by the definition of \(\lambda_{mn}\).
\end{proof}

\begin{remark}[Containment, not prime production]
The lattice theorem is one-way.  It places every realized key on the
displayed lattice; it does not assert that every integer \(z\) produces
prime largest factors or survives the largest-prime test, masks, complete
output labels, physical interval, or endpoints.
\end{remark}

\begin{corollary}[Physical primitive-row spacing]
\label{cor:primitive-row-spacing}
On the inherited physical carrier
\[
 (mn,h_0)=1,
\]
one has
\[
 (m,qu)=(n,qv)=1,\qquad
 \lambda_{mn}(q,u,v)=quv.
\]
Consequently,
\[
 j=j_0+quv\,z,\qquad
 Q_m=Q_{m,0}+mv\,z,\qquad
 Q_n=Q_{n,0}+nu\,z.
\]
\end{corollary}

\begin{proof}
Apply \eqref{eq:physical-row-primitivity} to both realized rows.
Since \(qu\mid mj+h_0\), any common divisor of \(m\) and \(qu\) divides
\(h_0\); hence it is \(1\).  The argument for \(n,qv\) is identical.
Now \(q,u,v\) are pairwise coprime, so
\(\operatorname{lcm}(qu,qv)=quv\).  Substitute this value into
\cref{eq:direct-time-ray,eq:direct-Qm-ray,eq:direct-Qn-ray}.
\end{proof}

\noindent
For a complete key \(i\), let \(\ell(i)\) denote the tuple of all its
labels other than the time coordinate.  This tuple retains, in
particular, the side, external block, opposite row, \(\gamma\)-coordinate,
orbit class, and every remaining orthogonal label.

\begin{corollary}[Ray census]
\label{cor:direct-ray-census}
If the allowed time interval has length \(J\), then
\[
 \#\{i\in\cI_{mn}:
       (q_i,u_i,v_i)=(q,u,v),\ \ell(i)=\ell_0\}
 \le
 1+\left\lfloor\frac{J}{\lambda_{mn}(q,u,v)}\right\rfloor
\]
for each fixed label tuple \(\ell_0\).  The unrestricted key count
requires summation over the realized \(\ell_0\)'s; no label multiplicity
is hidden in the displayed bound, and the sum does not exploit literal
phases.
\end{corollary}

\subsection{Frozen target sign}

\begin{theorem}[Primitive sign of one direct key]
\label{thm:direct-key-sign}
For every \(i\in\cI_{mn}\),
\begin{equation}
 \boxed{
 \mu(N_{i,m})\mu(N_{i,n})
 =
 \mu(u_i)\mu(v_i).}
 \label{eq:direct-primitive-sign}
\end{equation}
Consequently, the M\"obius sign is constant along every arithmetic ray in
\cref{thm:direct-entry-ray}, irrespective of whether its spacing
simplifies to \(quv\).
\end{theorem}

\begin{proof}
Each largest prime is coprime to its cofactor on the squarefree target
support, so
\[
 \mu(N_{i,m})=-\mu(S_{i,m}),\qquad
 \mu(N_{i,n})=-\mu(S_{i,n}).
\]
The two minus signs cancel.  The pairwise coprimality from
\cref{thm:canonical-primitive-key} gives
\[
 \mu(q_i u_i)\mu(q_i v_i)
 =
 \mu(q_i)^2\mu(u_i)\mu(v_i)
 =
 \mu(u_i)\mu(v_i).
\]
\end{proof}

\begin{corollary}[Same-residual sign check]
\label{cor:same-residual-sign-check}
At the distinguished point \((u,v)=(b,a)\),
\[
 \mu(u)\mu(v)=\mu(a)\mu(b)
 =\mu(d_m)\mu(d_n).
\]
This agrees with the same-residual sign-freezing theorem of TPC--68.
\end{corollary}

\begin{proof}
The first equality is symmetry.  Since the squarefree factors
\(g,a,b\) are pairwise coprime,
\[
 \mu(d_m)\mu(d_n)
 =\mu(ga)\mu(gb)
 =\mu(g)^2\mu(a)\mu(b).
\]
\end{proof}

\begin{remark}[What can oscillate]
The theorem freezes only the target M\"obius sign inside a fixed ray.
The literal coefficient
\(\overline{\beta_{i,m}}\beta_{i,n}\) may carry complex phases and
amplitude variation, and the set of surviving prime points may be highly
irregular.
\end{remark}
